\section{Physical metrics}

\subsection{Line-wise density statistics}

Chapter~3 highlighted the importance of the vertical organization of RTI
fields: the mixing region develops along the gravity direction, and the
position of the interface is one of the main physical signatures of the flow.
To preserve this information while reducing the dimensionality of the analysis, 
it is possible to compares the fields line by line (because of the horizontal periodicity).

Let $\rho_r(z,x)$ denote the $r$-th validation field and $\varrho_r(z,x)$ the
$r$-th field from a compared dataset, where $z$ is the vertical coordinate and
$x$ the horizontal coordinate. For each image and each horizontal line, the mean density
along the periodic direction is first computed as

\begin{equation}
    m_{\rho}(r,z)
    =
    \frac{1}{N_x}
    \sum_{x=1}^{N_x}\rho_r(z,x),
    \qquad
    m_{\varrho}(r,z)
    =
    \frac{1}{N_x}
    \sum_{x=1}^{N_x}\varrho_r(z,x).
\end{equation}

The ensemble mean profile of each line is then defined by

\begin{equation}
    \widetilde{\rho}_l(z)
    =
    \frac{1}{N_s}\sum_{r=1}^{N_s}m_{\rho}(r,z),
    \qquad
    \widetilde{\varrho}_l(z)
    =
    \frac{1}{N_s}\sum_{r=1}^{N_s}m_{\varrho}(r,z).
\end{equation}

This quantity represents the mean vertical density organization of the dataset.
It is directly related to the average position and thickness of the mixing
layer.

The ensemble variability of each line is described by the standard deviation

\begin{equation}
    \sigma_{\rho,l}(z)
    =
    \sqrt{
        \frac{1}{N_s}
        \sum_{r=1}^{N_s}
        \left(
            m_{\rho}(r,z)-\widetilde{\rho}_l(z)
        \right)^2
    },
\end{equation}

and similarly for $\sigma_{\varrho,l}(z)$. This profile measures how strongly
the mean value of a given line fluctuates across the dataset.

The distances used in the notebook are normalized $L^2$ distances between these
profiles:

\begin{equation}
    d_{\mu}(\rho,\varrho)
    =
    \frac{1}{N_z}
    \left\|
        \widetilde{\rho}_l
        -
        \widetilde{\varrho}_l
    \right\|_2,
    \qquad
    d_{\sigma}(\rho,\varrho)
    =
    \frac{1}{N_z}
    \left\|
        \sigma_{\rho,l}
        -
        \sigma_{\varrho,l}
    \right\|_2.
\end{equation}

These metrics are simple but physically informative: they test whether the
generated ensemble reproduces the average vertical density profile and the
amplitude of its fluctuations.

\begin{figure}[H]
    \centering
    \includegraphics[width=0.90\textwidth]{figures/ch4/skewnorm_joyplot.png}
    \caption{Line-wise density statistics extracted from the validation,
    generated and noise datasets. The skew-normal fits emphasize the vertical
    organization of the fields and the contrast between physical samples and
    pure noise.}
    \label{fig:ch4_line_statistics}
\end{figure}

\subsection{Fluctuation 'energy' spectrum}

The term 'energy spectrum' is used here in a loose sense to describe the distribution of fluctuations across spatial scales. It is not a true physical energy spectrum, but it has to be seen as signal energy; it would be equal to a real energy spectrum only if the density field were a velocity field. The first physical metric introduced above captures the mean vertical organization of the fields, but it does not provide any information about the distribution of fluctuations across scales. Starting from a dataset converted to a tensor of shape
$(r,z,x)$, the fluctuation field is defined in the notebook by subtracting the
global mean and taking its magnitude,

\begin{equation}
    \delta \rho_r(z,x)
    =
    \left|
        \rho_r(z,x) - \overline{\rho}
    \right|.
\end{equation}

A discrete Fourier transform is then applied, and the energy spectrum is
computed as

\begin{equation}
    E_{\rho}
    =
    \left|
        \mathcal{F}(\delta\rho)
    \right|^2.
\end{equation}

In the baseline, a one-dimensional transform is applied along one
spatial direction before reusing the same line-wise statistical machinery as
above. The resulting quantities therefore compare the average spectral content
and its variability across the dataset. Even though this construction is not a
full isotropic turbulence spectrum, it provides a first indication of whether
the generated fluctuations occupy approximately the correct range of spatial
scales.

\begin{figure}[H]
    \centering
    \includegraphics[width=0.90\textwidth]{figures/ch4/skewnorm_joyplot_energy_spectrum.png}
    \caption{Same comparison applied to the fluctuation energy spectrum. The
    generated samples remain much closer to the validation set than pure noise,
    but the spectral variability is still harder to match than the mean line
    profile.}
    \label{fig:ch4_energy_spectrum}
\end{figure}
